\documentclass[aspectratio=169,10pt]{beamer}
% Shared Beamer setup for Fall 2026 course slides.
% Clean Cal Poly-inspired theme: no custom class files or uncommon packages.
\mode<presentation>{
  \usetheme{default}
  \usefonttheme{professionalfonts}
}

\definecolor{CalPolyGreen}{HTML}{154734}
\definecolor{CalPolyGold}{HTML}{BD8B13}
\definecolor{CalPolySage}{HTML}{7FA28F}
\definecolor{CalPolyMint}{HTML}{DDF1E7}
\definecolor{CalPolySand}{HTML}{A29F86}
\definecolor{CalPolyBeige}{HTML}{EFEEDE}
\definecolor{DeckInk}{HTML}{1F2933}
\definecolor{DeckMuted}{HTML}{5F6B7A}
\definecolor{DeckSoft}{HTML}{F7F7F5}

\setbeamersize{text margin left=0.55in,text margin right=0.55in}
\setbeamercolor{normal text}{fg=DeckInk,bg=white}
\setbeamercolor{structure}{fg=CalPolyGreen}
\setbeamercolor{title}{fg=white}
\setbeamercolor{frametitle}{fg=CalPolyGreen,bg=white}
\setbeamercolor{block title}{bg=CalPolySage,fg=white}
\setbeamercolor{block body}{bg=CalPolyMint,fg=black}
\setbeamercolor{alerted text}{fg=CalPolyGold}

\setbeamerfont{title}{size=\Huge,series=\bfseries}
\setbeamerfont{frametitle}{size=\Large,series=\bfseries}
\setbeamerfont{block title}{size=\normalsize,series=\bfseries}

\setbeamertemplate{navigation symbols}{}
\setbeamertemplate{itemize item}{\color{CalPolyGold}$\blacktriangleright$}
\setbeamertemplate{itemize subitem}{\color{CalPolyGreen}$\bullet$}
\setbeamertemplate{footline}{
  \hfill{\color{DeckMuted}\scriptsize\insertframenumber/\inserttotalframenumber}
  \hspace{0.18in}\vspace{0.12in}
}
\setbeamertemplate{frametitle}{
  \vspace{0.18in}
  \hspace{0.02in}\insertframetitle\par
  \vspace{0.08in}
  \noindent\makebox[\linewidth][l]{%
    {\color{CalPolyGold}\rule{0.55in}{2pt}}%
    {\color{CalPolyGreen}\rule{\dimexpr\linewidth-0.55in\relax}{2pt}}%
  }
}

\newcommand{\CourseNumber}{Course}
\newcommand{\CourseTitle}{Course Title}
\newcommand{\LectureNumber}{0}
\newcommand{\LectureTitle}{Lecture Title}
\newcommand{\LectureDate}{Date}
\newcommand{\CourseUnit}{Unit}

\newcommand{\coursenumber}[1]{\renewcommand{\CourseNumber}{#1}}
\newcommand{\coursetitle}[1]{\renewcommand{\CourseTitle}{#1}}
\newcommand{\lecturenumber}[1]{\renewcommand{\LectureNumber}{#1}}
\newcommand{\lecturetitle}[1]{\renewcommand{\LectureTitle}{#1}}
\newcommand{\lecturedate}[1]{\renewcommand{\LectureDate}{#1}}
\newcommand{\courseunit}[1]{\renewcommand{\CourseUnit}{#1}}

\author{Kirk Duran}
\institute{California Polytechnic State University, San Luis Obispo}
\date{\LectureDate}

\newcommand{\makedecktitle}{
  \begingroup
  \setbeamercolor{background canvas}{bg=CalPolyGreen}
  \begin{frame}[plain]
    \vfill
    \begin{center}
      {\usebeamerfont{title}\color{white}\LectureTitle\par}
      \vspace{0.18in}
      {\Large\color{white}Lecture \LectureNumber\par}
    \end{center}
    \vfill
    \noindent{\color{white}\rule{\linewidth}{1.2pt}}\par
    \vspace{0.08in}
    \noindent\colorbox{CalPolyGold}{\makebox[0.24\linewidth][c]{\color{white}\Large\CourseNumber}}
    \hspace{0.12in}{\color{white}\Large Kirk Duran}
  \end{frame}
  \endgroup
}

\newenvironment{keyidea}{\begin{block}{Key idea}}{\end{block}}
\newenvironment{activity}{\begin{block}{In class}}{\end{block}}
\newenvironment{cpdefinition}{\begin{block}{Definition}}{\end{block}}
\newenvironment{exampleblockcp}{
  \begingroup
  \setbeamercolor{block title}{bg=CalPolySand,fg=white}
  \setbeamercolor{block body}{bg=CalPolyBeige,fg=black}
  \begin{block}{Example}
}{
  \end{block}
  \endgroup
}

\newcommand{\imageplaceholder}[2][0.3in]{%
  \noindent\makebox[\linewidth][c]{%
    \fcolorbox{CalPolySage}{CalPolyBeige}{%
      \parbox[c][#1][c]{0.86\linewidth}{%
        \centering
        {\color{DeckMuted}\scriptsize Image placeholder: }%
        {\color{DeckInk}\scriptsize\ttfamily\detokenize{#2}}%
      }%
    }%
  }%
}

\newcommand{\sectionframe}[1]{
  \begingroup
  \setbeamercolor{background canvas}{bg=CalPolyGreen}
  \begin{frame}[plain]
    \vfill
    {\color{white}\LARGE\bfseries #1\par}
    \vspace{0.18in}
    {\color{CalPolyGold}\rule{0.22\paperwidth}{3pt}}
    {\color{white}\rule{0.62\paperwidth}{3pt}}
    \vfill
  \end{frame}
  \endgroup
}

\newcommand{\placeholdernote}{
  \begin{block}{Placeholder}
    Replace these bullets with the final lecture material, examples, and in-class activities.
  \end{block}
}

\usepackage[T1]{fontenc}
\usepackage{lmodern}
\usepackage{amsmath,amssymb}
\usepackage{booktabs,array}
\pdfstringdefDisableCommands{\def\translate#1{#1}}
\graphicspath{{../assets/lecture-28/}}

% Lecture 28: Generalization and Ensemble Learning. Nov 2, 2026.
% Built to match the instructor's Lecture 26/27 Cal Poly Beamer conventions.
% Audience: senior-level AI without advanced mathematics prerequisites.
% Pedagogy: lead with meaning in ordinary English; calculations are supporting evidence.
% 50-minute plan: motivation 5, leakage/evaluation design 7, metrics 6,
% learning curves/model selection 7, ensemble/bagging 8, random forests 8,
% demonstrations 6, limitations 3. The broader ensemble survey is one slide.
% Python examples retain the instructor's original Wine and California Housing
% programs and add separate validation-aware random-forest comparisons.
% Original visualization source: supplied Lecture 19 Decision Trees, pp.10 and 12.
% Additional technical checks:
% https://scikit-learn.org/stable/modules/ensemble.html
% https://scikit-learn.org/stable/modules/generated/sklearn.ensemble.RandomForestClassifier.html
% https://scikit-learn.org/stable/modules/cross_validation.html
% https://scikit-learn.org/stable/modules/learning_curve.html
% The learning-curve and runtime diagrams are explicitly illustrative, not results.

\newcommand{\lecturevisual}[3][2.15in]{%
  \par\vspace{2pt}%
  {\centering\includegraphics[width=.98\linewidth,height=#1,keepaspectratio]{#2}\par}%
  {\centering\scriptsize\color{DeckMuted}#3\par}%
  \vspace{2pt}%
}
\newcommand{\conceptlead}[1]{%
  {\small\bfseries\color{CalPolyGreen}#1\par}%
}
\newenvironment{calculation}[1][How we calculate it]{%
  \par\begingroup
  \setbeamercolor{calculation body}{bg=DeckSoft,fg=DeckInk}%
  \begin{beamercolorbox}[sep=4pt,wd=\linewidth]{calculation body}%
  \footnotesize{\bfseries\color{DeckMuted}#1\par}%
  \setlength{\abovedisplayskip}{3pt}%
  \setlength{\belowdisplayskip}{3pt}%
  \setlength{\jot}{0pt}%
}{\end{beamercolorbox}\endgroup}
\setbeamersize{text margin left=.42in,text margin right=.42in}
\coursenumber{CS 4880}
\coursetitle{Artificial Intelligence}
\lecturenumber{28}
\lecturetitle{Generalization and Ensemble Learning}
\lecturedate{November 2, 2026}
\courseunit{Learning Agents and Machine Learning}

\begin{document}
\makedecktitle

\section{Why a Perfect Tree May Fail}

\begin{frame}[t]{Why Can a Perfect Decision Tree Be a Bad Model?}
\small
\conceptlead{Memorizing old outcomes and predicting future outcomes are different achievements.}
\begin{columns}[T,onlytextwidth]
 \begin{column}{.52\textwidth}
 \begin{itemize}
  \item Last lecture, an unrestricted tree could split until nearly every training example was classified correctly.
  \item A rule like ``this exact camera pattern means blocked'' may describe \emph{one recording}, not a useful property of corridors.
  \item The robot must succeed on \textbf{new delivery runs}, new buildings, and new conditions.
 \end{itemize}
 \begin{keyidea}The goal is \textbf{generalization}: reliable predictions on observations not used to fit or choose the model.\end{keyidea}
 \end{column}
 \begin{column}{.44\textwidth}
 \begin{block}{Two different reports}
 \textbf{Training accuracy:} 100\%\\[4pt]
 \textbf{New-run accuracy:} 67\%\\[7pt]
 A high first number does not guarantee a high second one.\par\smallskip
 \scriptsize Hypothetical example, not a measured result.
 \end{block}
 \end{column}
\end{columns}
\end{frame}

\begin{frame}[t]{Three Ways Predictions Fail}
\small
\conceptlead{Ask whether the model is too simple, too specific, or facing a changed world.}
\lecturevisual[1.58in]{three_failures.png}{}
\begin{itemize}
 \item \textbf{Underfitting:} a depth-one tree misses important interactions.
 \item \textbf{Overfitting:} a deep tree turns unusual training examples into fragile rules.
 \item \textbf{Distribution shift:} the deployment environment differs from training (e.g., rain, lighting, or a new campus building).
\end{itemize}
\end{frame}

\section{Design an Honest Evaluation}

\begin{frame}[t]{Train, Validate, Test: Each Set Has One Job}
\small
\conceptlead{Keep model fitting, model selection, and final reporting separate.}
\begin{center}
\renewcommand{\arraystretch}{1.17}
\begin{tabular}{>{\raggedright\arraybackslash}p{.19\linewidth}>{\raggedright\arraybackslash}p{.33\linewidth}>{\raggedright\arraybackslash}p{.36\linewidth}}
\toprule
\textbf{Partition} & \textbf{What we do with it} & \textbf{Question it answers}\\
\midrule
\textbf{Training} & Fit weights, splits, and leaf predictions. & What relationships can the model learn?\\
\textbf{Validation} & Choose depth, forest size, and threshold. & Which candidate looks best on unseen cases?\\
\textbf{Test} & Evaluate the frozen selected model \emph{once}. & How well did the chosen procedure transfer?\\
\bottomrule
\end{tabular}
\end{center}
\begin{keyidea}Set hyperparameters on validation, learn splits on training, and report the final test only after choosing the model.\end{keyidea}
\end{frame}

\begin{frame}[t]{Data Leakage: Different Frames, Same Experience}
\small
\conceptlead{A random row split is not independent if rows record nearly the same event.}
\lecturevisual[1.16in]{leakage_groups.png}{Near-duplicate recording frames can leak across a frame-level split.}
\begin{itemize}
 \item Adjacent frames share nearly the same corridor, lighting, and obstacles.
 \item Sharing a run between training and test gives an \emph{overly optimistic} score.
 \item Keep delivery runs/sessions intact. To predict future runs, hold out later time periods; random grouped splits alone do not preserve time order.
 \item Fit preprocessing (e.g., imputation) on training data only, inside each fold.
\end{itemize}
\end{frame}

\begin{frame}[t]{Cross-Validation: Do Not Trust One Lucky Split}
\small
\conceptlead{Rotate held-out groups, refit, and summarize the validation scores.}
\begin{columns}[T,onlytextwidth]
\begin{column}{.53\textwidth}
\lecturevisual[1.13in]{group_folds.png}{T = train; V = validate; test stays untouched.}
\end{column}
\begin{column}{.44\textwidth}
\begin{itemize}
 \item $K$ folds: train $K$ times, validating on a different fold each time.
 \item Grouped folds keep each delivery run intact.
 \item The spread measures split sensitivity, not a confidence interval.
\end{itemize}
\end{column}
\end{columns}
\begin{calculation}[Average the held-out scores]
$\displaystyle\overline s=\frac1K\sum_{k=1}^K s_k$ is the mean held-out score.\par
$K$: fold count; $k$: fold index; $s_k$: held-out score; $\overline s$: mean score. Refit preprocessing within each training fold; leave the test untouched.
\end{calculation}
\end{frame}

\section{What Counts as a Good Prediction?}

\begin{frame}[t]{A Confusion Matrix Names the Mistakes}
\small
\conceptlead{On 40 hypothetical segments, count the different mistake types.}
\begin{columns}[T,onlytextwidth]
 \begin{column}{.50\textwidth}
 \begin{center}
 \renewcommand{\arraystretch}{1.2}
 \begin{tabular}{lcc}
  & \textbf{Actual} & \textbf{Actual}\\[-4pt]
  & \textbf{Blocked} & \textbf{Clear}\\\hline
 \textbf{Pred. Blocked} & TP = 9 & FP = 3\\
 \textbf{Pred. Clear} & FN = 6 & TN = 22\\
 \end{tabular}
 \end{center}
 \textbf{False positive:} unnecessary rerouting.\par\smallskip
 \textbf{False negative:} the robot enters an unexpectedly blocked route.\par\smallskip
 High accuracy can hide missed blockages: inspect false negatives and recall.

 \end{column}
 \begin{column}{.47\textwidth}
 \begin{calculation}[Four ways to describe performance]
 \[\begin{aligned}
 \text{Accuracy}&=\frac{TP+TN}{N}=\frac{31}{40}=77.5\%\\
 \text{Precision}&=\frac{TP}{TP+FP}=\frac9{12}=75\%\\
 \text{Recall}&=\frac{TP}{TP+FN}=\frac9{15}=60\%\\
 F_1&=\frac{2TP}{2TP+FP+FN}=\frac{18}{27}\approx66.7\%
 \end{aligned}\]
 \scriptsize Positive = Blocked. TP/TN: true positives/negatives; FP/FN: false positives/negatives. $N=40$ cases. $F_1$: harmonic mean of precision and recall.
 \end{calculation}
 \end{column}
\end{columns}

\end{frame}

\begin{frame}[t]{A Decision Threshold Turns Scores Into Labels}
\small
\conceptlead{The same learned score can produce different labels under different risk tolerances.}
\begin{columns}[T,onlytextwidth]
 \begin{column}{.50\textwidth}
 \begin{itemize}
  \item Suppose a model assigns blockage score $p=0.70$ to a corridor.
  \item Threshold $0.50$: predict \textbf{blocked}; threshold $0.80$: predict \textbf{clear}.
  \item Lowering the threshold cannot reduce recall on a fixed dataset, but adds predicted positives. Precision may rise or fall.
 \end{itemize}
 \end{column}
 \begin{column}{.47\textwidth}
 \lecturevisual[1.73in]{threshold_tradeoff.png}{Illustrative precision--recall tradeoff when a threshold changes.}
 \end{column}
\end{columns}

\end{frame}

\begin{frame}[t]{Regression Error: How Far Off Are the Predictions?}
\small
\conceptlead{Numerical predictions need an error magnitude, not a right/wrong count.}
\begin{calculation}[Regression error, if the target is delay in minutes]
\[\operatorname{MAE}=\frac1m\sum_i|\widehat y_i-y_i|,\qquad
\operatorname{RMSE}=\sqrt{\frac1m\sum_i(\widehat y_i-y_i)^2}.\]
$m$: evaluated cases; $i$: case index; $y_i$: true delay; $\widehat y_i$: predicted delay. MAE is mean absolute error; RMSE is root mean squared error. Both use minutes; RMSE emphasizes large misses.
\end{calculation}
\begin{block}{A two-case example}
True delays: $[2,6]$ minutes; predictions: $[3,3]$ minutes. Absolute errors are $[1,3]$: MAE $=2$ minutes and RMSE $=\sqrt{5}\approx2.24$ minutes.
\end{block}
\begin{keyidea}A larger miss influences RMSE more strongly because errors are squared before averaging.\end{keyidea}
\end{frame}

\section{Diagnose and Select Models}

\begin{frame}[t]{Learning Curves: Does More Training Data Help?}
\small
\conceptlead{A learning curve shows training and validation performance as sample size grows.}
\begin{columns}[T,onlytextwidth]
 \begin{column}{.55\textwidth}
 \lecturevisual[2.08in]{learning_curve.png}{}
 \end{column}
 \begin{column}{.41\textwidth}
 \begin{itemize}
 \item \textbf{Both errors high:} likely underfitting or weak features.
 \item \textbf{Training error low, validation error high:} likely overfitting.
 \item More representative data may narrow the gap; improvement is not guaranteed.
 \end{itemize}
 \end{column}
\end{columns}
\end{frame}

\begin{frame}[t]{Choose the Candidate That Works on Validation}
\small
\conceptlead{The strongest training score need not belong to the most useful model.}
\begin{columns}[T,onlytextwidth]
 \begin{column}{.51\textwidth}
 \lecturevisual[1.98in]{model_validation.png}{}
 \end{column}
 \begin{column}{.46\textwidth}
 \begin{itemize}
  \item The unrestricted tree gets \textbf{100\% training accuracy}, but only 72\% on validation.
  \item The random forest gets the best validation score here: \textbf{88\%}.
  \item Choose hyperparameters using validation; evaluate the final selected model on test data once.
 \end{itemize}
 \begin{keyidea}This example illustrates selection, not a guarantee that forests always win.\end{keyidea}
 \end{column}
\end{columns}
\end{frame}

\begin{frame}[t]{Why More Data Helps --- and What It Cannot Fix}
\small
\conceptlead{Additional independent experience makes sample-specific accidents less persuasive.}
\begin{itemize}
 \item More \emph{representative} delivery runs show the model what varies and which patterns persist.
 \item Extra copies of nearly identical frames do not add the same amount of new information.
 \item More of the \emph{same} data cannot fix missing sensor features, mislabeled outcomes, or a new building unlike the training environments.
 \item If one validation split gives an unusually good or bad result, grouped cross-validation can expose that instability.
\end{itemize}
\begin{block}{Transition to ensembles}
Rather than trust one tree trained on one sample, could we combine many trees that make \textbf{different mistakes}?
\end{block}
\end{frame}

\section{Bagging: Many Different Trees}

\begin{frame}[t]{Ensemble Learning: Combine Imperfect Predictors}
\small
\conceptlead{One uncertain judgment can become more stable when several useful judgments are combined.}
\begin{itemize}
 \item An \textbf{ensemble} uses multiple learned models to produce one prediction.
 \item For \textbf{classification}, models can vote on the label or average their class probabilities.
 \item For \textbf{regression}, models can average their numerical predictions.
 \item Shared mistakes persist; diverse errors are more likely to cancel.
\end{itemize}
\begin{calculation}[A regression ensemble in compact form]
\[\widehat y_{\rm ensemble}(\mathbf x)=\frac1B\sum_{b=1}^B\widehat y_b(\mathbf x).\]
$\mathbf x$: feature vector for one case; $B$: model count; $b$: model index; $\widehat y_b(\mathbf x)$: prediction of model $b$. Average these to get the ensemble prediction.
\end{calculation}
\end{frame}

\begin{frame}[t]{Bagging: Bootstrap Aggregation}
\small
\conceptlead{Train each tree on a newly sampled version of the same training set.}
\lecturevisual[1.02in]{bagging_pipeline.png}{Bootstrap samples train separate trees, whose outputs are combined.}
\begin{enumerate}
 \item Draw $N$ training rows \textbf{with replacement}; repetition is allowed.
 \item Fit a separate tree to that sample. Repeat for many bootstrap samples.
 \item Combine results by majority vote or averaging.
\end{enumerate}
\begin{calculation}[Why the samples really differ]
$N$: original training-set size and draws per sample. The expected distinct fraction is $1-(1-1/N)^N$, about $63\%$ for large $N$. Related frames should not be mistaken for independent evidence.
\end{calculation}
\end{frame}

\begin{frame}[t]{Five Trees Vote: A Worked Prediction}
\small
\conceptlead{The ensemble chooses the class with the most votes --- not necessarily a unanimous answer.}
\lecturevisual[0.98in]{vote_five.png}{Three Blocked votes and two Clear votes.}
\begin{calculation}[Classification vote]
\[\widehat y(\mathbf x)=\operatorname{mode}\!\left(\widehat y_1(\mathbf x),\ldots,\widehat y_B(\mathbf x)\right).\]
$\mathbf x$: one case; $\widehat y_b$: label from tree $b$; $B$: tree count. \emph{Mode} is the most frequent label: Blocked wins $3$ to $2$ here. Specify a tie rule if needed.
\end{calculation}
\smallskip\scriptsize A 3-to-2 vote is not a calibrated 60\% probability. In scikit-learn, random forests average tree probabilities rather than hard-vote; shared mistakes can persist.
\end{frame}

\section{Random Forests}

\begin{frame}[t]{Why Bagging Similar Trees Helps Less}
\small
\conceptlead{Correlated errors survive a vote: the trees need reasons to disagree.}
\lecturevisual[0.88in]{correlated_errors.png}{}
\smallskip\begin{itemize}
 \item Bootstrap samples vary, but dominant features may repeat. If trees fail on the \emph{same cases}, averaging helps less.
\end{itemize}
\begin{calculation}[Why correlation matters (optional intuition)]
\[\operatorname{Var}(\text{average})=\sigma^2\left(\rho+\frac{1-\rho}{B}\right).\]
\scriptsize $B$: model count; $\sigma^2$: variance of each estimate across training samples; $\rho$: common pairwise correlation. $\rho=1$: no variance reduction; $\rho=0$: variance divided by $B$. This is a variance model, not an accuracy formula.
\end{calculation}
\end{frame}

\begin{frame}[t]{Random Forests Add a Second Source of Randomness}
\small
\conceptlead{Random forests randomize the rows used per tree and the candidate features considered at each split.}
\lecturevisual[1.36in]{forest_two_randomness.png}{Randomized observations and random feature subsets encourage diverse trees.}
\begin{itemize}
 \item \textbf{Bootstrap rows:} different training samples for different trees.
 \item \textbf{Random feature candidates:} each split searches only a subset of the inputs.
 \item Each tree chooses its best available split; feature randomness makes trees less alike.
\end{itemize}
\end{frame}

\begin{frame}[t]{How a Random Forest Trains and Predicts}
\small
\conceptlead{Many ordinary decision trees are trained; the final output combines their answers.}
\begin{columns}[T,onlytextwidth]
 \begin{column}{.52\textwidth}
 \begin{enumerate}
  \item Create a bootstrap sample for each tree.
  \item Recursively split with a random subset of candidate features at each node.
  \item Let all trees predict the new observation.
  \item \textbf{Classify:} average class probabilities. \textbf{Regress:} average numerical values.
 \end{enumerate}
 \end{column}
 \begin{column}{.44\textwidth}
 \begin{block}{Important hyperparameters}
 \texttt{n\_estimators}: tree count.\\[4pt]
 \texttt{max\_depth}: depth limit.\\[4pt]
 \texttt{min\_samples\_leaf}: smallest leaf.\\[4pt]
 \texttt{max\_features}: features considered per split.
 \end{block}
 \end{column}
\end{columns}
\begin{calculation}[One common feature-sampling convention]
With $d$ input features, \texttt{max\_features="sqrt"} considers about $\sqrt d$ candidate features per split ($\sqrt{\cdot}$: square root). Regression defaults may use all features; the setting is tunable.
\end{calculation}
\begin{keyidea}More trees may stabilize predictions; gains diminish while compute cost grows.\end{keyidea}
\end{frame}

\begin{frame}[t]{Beyond Random Forests: Other Ensemble Methods}
\small
\conceptlead{Ensemble learning is a family of approaches, not one algorithm.}
\lecturevisual[1.02in]{ensemble_families.png}{Parallel bagging, sequential boosting, and learned stacking illustrated.}
\vspace{-2pt}
\renewcommand{\arraystretch}{1.10}
\begin{center}
\scriptsize
\begin{tabular}{p{.18\linewidth}p{.68\linewidth}}
\textbf{Method} & \textbf{How it combines models}\\\hline
\textbf{Bagging} & Separately fitted bootstrap models combined by averaging/voting; aims to reduce variance.\\
\textbf{Random forest} & Bagged trees plus random feature candidates to reduce tree correlation.\\
\textbf{Boosting} & Sequentially adds models to improve errors/loss (AdaBoost, gradient boosting, XGBoost).\\
\textbf{Stacking} & A meta-model learns from out-of-fold base-model predictions to avoid training leakage.\\
\textbf{Voting} & Independently trained classifiers combine hard labels or probabilities.\\
\end{tabular}
\end{center}
\end{frame}

\section{Python: From Your Original Trees to a Forest}

\begin{frame}[fragile,t]{Python Example 1: Explore the Wine Tree}
\small
\conceptlead{A tree fitted on all data illustrates rules; it does not measure generalization.}
\begin{columns}[T,onlytextwidth]
\begin{column}{.40\textwidth}
\lecturevisual[1.70in]{wine_tree.png}{Top levels of the depth-4 Wine tree from Lecture 27; fitted on all data.}
\end{column}
\begin{column}{.57\textwidth}
\scriptsize
\begin{verbatim}
import matplotlib.pyplot as plt
from sklearn.datasets import load_wine
from sklearn.tree import (
    DecisionTreeClassifier, plot_tree)
wine = load_wine()
X, y = wine.data, wine.target
clf = DecisionTreeClassifier(max_depth=4, random_state=42)
clf.fit(X, y)
plt.figure(figsize=(14, 7))
plot_tree(clf, max_depth=2, filled=True,
          feature_names=wine.feature_names,
          class_names=wine.target_names, fontsize=8)
plt.show()
\end{verbatim}
\end{column}
\end{columns}
\end{frame}

\begin{frame}[fragile,t]{Wine: Compare on Validation, Then Report on Test}
\small
\conceptlead{Use the same partitions for both candidates; keep the final test out of selection.}
\begin{columns}[T,onlytextwidth]
\begin{column}{.49\textwidth}
\scriptsize
\begin{verbatim}
# Continue after loading Wine above.
from sklearn.model_selection import train_test_split
from sklearn.ensemble import RandomForestClassifier
X_tr, X_hold, y_tr, y_hold = train_test_split(
    X, y, test_size=.4, stratify=y,
    random_state=42)
X_v, X_test, y_v, y_test = train_test_split(
    X_hold, y_hold, test_size=.5,
    stratify=y_hold, random_state=42)
\end{verbatim}
\end{column}
\begin{column}{.49\textwidth}
\scriptsize
\begin{verbatim}
models = [
    DecisionTreeClassifier(
        max_depth=4, random_state=42),
    RandomForestClassifier(n_estimators=100,
        max_depth=6, random_state=42)]
for model in models:
    model.fit(X_tr, y_tr)
scores = [m.score(X_v, y_v) for m in models]
best = models[scores.index(max(scores))]
print("Validation accuracy:", scores)
print("Test:", best.score(X_test, y_test))
\end{verbatim}
\end{column}
\end{columns}
\begin{calculation}[Read the variable names]
\texttt{tr}: training; \texttt{v}: validation; \texttt{hold}: temporary held-out pool. \texttt{score} is accuracy here; select its maximum. This small independent-sample teaching example uses stratification, not grouped deployment evaluation.
\end{calculation}
\end{frame}

\begin{frame}[fragile,t]{Python Example 2: Explore a Housing Regression Tree}
\small
\conceptlead{A regression tree predicts a leaf's mean target for the squared-error criterion.}
\begin{columns}[T,onlytextwidth]
\begin{column}{.40\textwidth}
\lecturevisual[1.70in]{california_tree.png}{Top levels of the depth-4 California Housing tree from Lecture 27.}
\end{column}
\begin{column}{.57\textwidth}
\scriptsize
\begin{verbatim}
from sklearn.datasets import (
    fetch_california_housing)
from sklearn.tree import DecisionTreeRegressor
data = fetch_california_housing()
X, y = data.data, data.target
# Structural illustration, not held-out evaluation.
tree = DecisionTreeRegressor(max_depth=4, random_state=42)
tree.fit(X, y)
plt.figure(figsize=(14, 7))
plot_tree(tree, max_depth=2, filled=True,
          feature_names=data.feature_names, fontsize=8)
plt.show()
\end{verbatim}
\end{column}
\end{columns}
\scriptsize The target is district median house value, in units of \$100,000; rows describe census block groups, not individual houses.
\end{frame}

\begin{frame}[fragile,t]{Housing: Select by Validation Error, Then Test Once}
\small
\conceptlead{For regression, average the tree predictions and prefer lower validation error.}
\begin{columns}[T,onlytextwidth]
\begin{column}{.49\textwidth}
\scriptsize
\begin{verbatim}
# Continue with Housing X, y.
from sklearn.ensemble import RandomForestRegressor
from sklearn.metrics import mean_absolute_error
from sklearn.metrics import root_mean_squared_error
X_tr, X_hold, y_tr, y_hold = train_test_split(
    X, y, test_size=.4, random_state=42)
X_v, X_test, y_v, y_test = train_test_split(
    X_hold, y_hold, test_size=.5, random_state=42)
models = [
    DecisionTreeRegressor(
        max_depth=4, random_state=42),
    RandomForestRegressor(n_estimators=100,
        max_depth=12, random_state=42, n_jobs=-1)]
\end{verbatim}
\end{column}
\begin{column}{.49\textwidth}
\scriptsize
\begin{verbatim}
for model in models:
    model.fit(X_tr, y_tr)
errors = [mean_absolute_error(y_v, m.predict(X_v))
          for m in models]
best = models[errors.index(min(errors))]
pred = best.predict(X_test)
print("Validation MAE:", errors)
print("Test MAE:",
      mean_absolute_error(y_test, pred))
print("Test RMSE:",
      root_mean_squared_error(y_test, pred))
\end{verbatim}
\end{column}
\end{columns}
\begin{calculation}[What this simple split does and does not test]
Both metrics use \$100,000 units. Random splits demonstrate the workflow; geographic holdouts are needed to assess transfer to new regions. Code assumes scikit-learn 1.4+ for \texttt{root\_mean\_squared\_error}.
\end{calculation}
\end{frame}

\begin{frame}[t]{Fair Comparisons: Same Inputs, Same Protocol}
\small
\conceptlead{Changing the model while holding evaluation fixed turns a demonstration into evidence.}
\begin{columns}[T,onlytextwidth]
 \begin{column}{.54\textwidth}
 \lecturevisual[1.78in]{controlled_protocol.png}{Grouped partitions determine when a model may see each observation.}
 \end{column}
 \begin{column}{.43\textwidth}
 \begin{itemize}
 \item Fit a \textbf{linear baseline}, shallow tree, unrestricted tree, and forest on the same features and train split.
 \item Tune on validation folds; record depth, accuracy, and runtime.
 \item Freeze the selected model, then evaluate on test once.
 \end{itemize}
 \end{column}
\end{columns}
\begin{calculation}[What to vary in the lab]
\texttt{n\_estimators} changes forest size; \texttt{max\_depth} changes each tree's complexity. Report validation metrics and training-time cost together.
\end{calculation}
\end{frame}

\section{Limits and Synthesis}

\begin{frame}[t]{More Trees Are Not a Free Improvement}
\small
\conceptlead{Forests can be more stable than one deep tree, but they cost computation and cannot repair a broken evaluation.}
\begin{columns}[T,onlytextwidth]
 \begin{column}{.53\textwidth}
 \lecturevisual[1.76in]{forest_runtime.png}{Illustrative learning-performance saturation and training-time tradeoff.}
 \end{column}
 \begin{column}{.44\textwidth}
 \begin{itemize}
 \item Gains can level off while training time continues to grow.
 \item Highly correlated or systematically biased trees may all make the same mistake.
 \item A forest is generally less interpretable than one small decision tree.
 \item No ensemble automatically fixes leakage, missing sensors, or distribution shift.
 \end{itemize}
 \end{column}
\end{columns}
\end{frame}

\begin{frame}[t]{Synthesis: Better Models Require Better Evidence}
\small
\begin{enumerate}
 \item \textbf{Generalization:} evaluate future performance, not just how well the model fits old examples.
 \item \textbf{Reliable evaluation:} prevent leakage; group correlated observations; use validation to choose settings and test once.
 \item \textbf{Error diagnosis:} confusion matrices, regression error, and learning curves reveal different kinds of failure.
 \item \textbf{Ensemble learning:} bagging combines bootstrap-fitted predictors; random forests add feature randomness to lower correlation.
 \item \textbf{Limits:} averaging can reduce variance, not recover information that was never measured.
\end{enumerate}
\begin{keyidea}When would five identical trees fail together, and why would five diverse trees do better?\end{keyidea}
\end{frame}

\end{document}
